\clearpage
\fancyhead[L]{\small CONNAISSANCE DES TEMS FÜR 1818.}
\hypertarget{native-600-lower}{}
\newcommand{\ObservedJupiterErrorMark}{\hbox{\vrule width0pt height4.000000bp depth0pt\special{pdf:bcontent}\special{pdf:code q 0 g 0.500000 4.000000 m 0.000000 3.500000 l 0.000000 3.250000 l 0.750000 2.500000 l 0.750000 2.250000 l 1.000000 2.000000 l 1.000000 1.750000 l 0.750000 1.500000 l 0.250000 1.500000 l 0.000000 1.250000 l 0.000000 0.500000 l 0.250000 0.250000 l 1.250000 0.250000 l 1.500000 0.000000 l 2.000000 0.000000 l 2.250000 0.250000 l 2.500000 0.250000 l 3.250000 1.000000 l 3.250000 1.250000 l 2.500000 2.000000 l 2.500000 2.500000 l 3.000000 3.000000 l 3.000000 3.250000 l 2.250000 4.000000 l 1.750000 4.000000 l 1.500000 3.750000 l 1.250000 3.750000 l 1.000000 4.000000 l h f* Q}\special{pdf:econtent}\kern3.500000bp}}

\begin{center}\small Göttingische gelehrte Anzeigen. Stück 54. Seite 531..536. 1817 April 5.\\[3pt]\rule{.72\linewidth}{.3pt}\end{center}

\emph{Connaissance des tems, ou des mouvemens célestes, à l'usage des astronomes et des navigateurs, pour l'an \textup{1818}; publiée par le bureau des longitudes.} Paris 1815. Bei der Witwe Courcier. 412 Seiten in Octav.

Die Einrichtung des astronomischen Kalenders hat in diesem Jahrgange einige kleine Veränderungen erlitten. Die gerade Aufsteigung des Mondes, welche sonst nur in Bogenminuten angesetzt wurde, ist jetzt für Mittag und Mitternacht auf Secunden berechnet, wogegen die Declinationen, die sonst von 6 zu 6 Stunden angegeben wurden, jetzt, in Minuten, nur für Mittag und Mitternacht angesetzt sind. Diese Abänderung muss denjenigen Astronomen angenehm sein, welche durch correspondirende Beobachtungen von Mondsculminationen den Unterschied ihrer geographischen Längen bestimmen wollen. Ferner sind die einzelnen Phänomene, welche sonst bei jedem Monat besonders aufgeführt wurden, jetzt zweckmässiger am Ende des Kalenders zusammengestellt. Der dadurch gewonnene Platz ist zur Angabe der heliocentrischen Längen und Breiten der Planeten, und ihrer geraden Aufsteigungen in Zeit verwandt. Es ist nur zu bedauern, dass man letztere blos auf Minuten (also Viertelsgrade in Bogen) angesetzt hat; es ist sehr oft wünschenswerth, um die Stellung der Planeten gegen benachbarte Fixsterne im voraus beurtheilen zu können, ihren Platz auf die Bogenminute genau vorher zu wissen, und anstatt der heliocentrischen Oerten, welche in der Ephemeride wenig oder gar keinen Nutzen haben, wäre die Ansetzung des Abstandes von der Erde bei weitem wichtiger. Es ist sehr zu bedauern, dass keine einzige unsrer astronomischen Ephemeriden dies Element mittheilt, obgleich dasselbe für Aberration, Parallaxe und Reduction der gemessenen Durchmesser gleich wichtig ist.

Unter den beigefügten Aufsätzen machen einige Auszüge, welche Herr \textsc{Delambre} aus verschiedenen Englischen Schriften gemacht hat, den Anfang. Da diese auf dem festen Lande doch noch immer in wenige Hände kommen, so werden diese Auszüge manchem Leser angenehm sein: auch die beigefügten Anmerkungen von einem so erfahrnen Praktiker liest man gern, wenn gleich derselbe sich oft wiederholt. Mit ermüdender und unnöthiger Weitschweifigkeit ist dagegen der folgende Aufsatz geschrieben, über eine Aufgabe von \textsc{Regiomontan}, die Umstände anzugeben, unter denen zwei Punkte der Ekliptik in Länge eben so viel verschieden sind, wie in gerader Aufsteigung: diese unbedeutende und \SourceLeaf{601}leichte Aufgabe hätte in der Connoissance des tems keine 10 Seiten verdient. --- Neue Tafeln für die Aberration der Planeten in Länge und Breite von \textsc{Puissant}. Diese Tafeln haben eine so bequeme Einrichtung, wie es die Natur des Gegenstandes erlaubt, und mögen in dem Falle mit Nutzen angewandt werden, wo die Entfernung des Planeten von der Erde ganz unbekannt ist. Bei den Formeln, wonach die Tafeln berechnet sind, bemerkt der Verfasser: Ces formules sont \emph{complètes}; il est remarquable en outre, qu'elles sont exactes, \emph{aux quantités près du troisième ordre}; car les termes en $\mathrm{e}^2$ se detruisent mutuellement, comme je m'en suis assuré. Vermuthlich soll diese etwas dunkel ausgedrückte Stelle bedeuten, dass der Verf. die Aberration nach den Potenzen der Excentricität entwickelt hat, und dass die von der Excentricität unabhängigen Glieder, eben so wie die von der ersten Ordnung der Excentricität vollständig sind. Ist diese Auslegung (ohne welche die unterstrichenen Ausdrücke im Widerspruch zu stehen scheinen) die richtige, so scheint dem Verf. entgangen zu sein, dass seine Formeln in der That nicht blos bis zu den Grössen der dritten Ordnung der Excentricität genau sind, sondern \emph{absolute} Vollständigkeit haben. --- Von demselben Verf. eine Anmerkung zu einem Aufsatze von \textsc{Lagrange}, die Parallaxenrechnung bei den Finsternissen betreffend, wo Herr \textsc{Puissant} den von \textsc{Lagrange} für die scheinbare Entfernung der Mittelpunkte der Sonne und des Mondes gegebenen Formeln durch Einführung von Hülfsmitteln eine für die numerische Rechnung geschmeidigere Gestalt zu geben sucht: bequemer ist jedoch hier die Anwendung der Tafel, aus der man sofort die Logarithmen der Summen und Differenzen von Grössen findet, die unmittelbar durch ihre Logarithmen gegeben sind. --- Die neue allgemeine Tafel für die parabolische Bewegung der Cometen von \textsc{Burckhardt} verdient den Vorzug vor der gewöhnlichen Tafel sowohl, als auch selbst vor der \textsc{Barker}'schen; sie unterscheidet sich von jener, dass nicht die Zeit, die seit dem Durchgange durch das Perihelium verflossen ist, sondern der Logarithme dieser Zeit ihr Argument ist. Verglichen mit der gewöhnlichen Tafel erspart sie das Aufschlagen eines Logarithmen; eben den Vorzug hat sie vor der \textsc{Barker}'schen Tafel bis zu $45^\circ$ wahrer Anomalie, auch ist das Interpoliren in der \textsc{Burckhardt}'schen Tafel meistens etwas bequemer. --- Neue Bestimmung der Bahn der Vesta von \textsc{Daussy}. Herr \textsc{Daussy} hat sich durch die Berechnung der Störungen der Bewegung der Vesta durch Jupiter, Saturn und Mars nach der \textsc{Laplace}'schen Methode, welche bei der mässigen Excentricität und Neigung jenes Planeten als zulänglich betrachtet werden kann, sehr verdient gemacht, und die ersten sieben Oppositionen zeigen in der That eine sehr befriedigende Uebereinstimmung. --- Ueber einen neuen Apparat zur Vergleichung linearischer Masse von \textsc{Prony}. Das vornehmste Stück dieses Apparats ist ein auf Glas in 100 Theile eingetheilter Millimeter; mehrere Französische Künstler liefern diese Theilung mit einer in der That bewundernswürdigen Feinheit und Genauigkeit. Von den zu vergleichenden Massen müssen solche Zusammensetzungen gemacht werden, dass ihr Unterschied höchstens ein Paar Millimeter beträgt, und blos dieser Unterschied wird durch den Apparat bestimmt, dessen vollständigere Beschreibung hier zu weitläuftig sein würde. --- Ueber die Ebbe und Fluth des Meeres von \textsc{Laplace}. Dieser kleine Aufsatz ist die Einleitung zu neuen interessanten Untersuchungen, welche dieser grosse Geometer über die zu Brest angestellten Beobachtungen der Ebbe und Fluth gemacht hat, und muss auf das nähere Detail dieser Untersuchungen sehr begierig machen. Besonders merkwürdig ist die daraus abgeleitete Bestimmung der Mondsmasse $(= \frac{1}{68.7})$, und der Nutationsconstante $(= 9''65)$, um so mehr, da die von Hrn. \textsc{von Lindenau} mit so grosser Sorgfalt discutirten Beobachtungen des Polarsterns einen beträchtlich kleinern Werth der letztern gegeben haben. --- Die beiden folgenden Aufsätze von demselben Verfasser über die Anwendung der Wahrscheinlichkeitsrechnung auf die Naturwissenschaft beschäftigen sich mit der Theorie der Genauigkeit der nach der Methode der kleinsten Quadrat gefundenen Resultate, und der Bestimmung der Genauigkeit der Beobachtungen selbst. Die Resultate dieser Untersuchung kommen im Grunde mit dem, was von \textsc{Gauss} in der Theoria Motus \SourceLeaf{602}Corporum Coelestium, und in einem Aufsatze in der Zeitschrift für Astronomie entwickelt ist, ganz überein, obgleich \textsc{Laplace} den Gegenstand aus einem etwas verschiedenen Gesichtspunkte betrachtet hat. Es werden hier zugleich mehrere interessante Anwendungen dieser Theorie gegeben. Eine davon ist die Bestimmung der Jupitersmasse aus den Störungen, welche Saturn durch die Einwirkung des Jupiter erleidet. Das aus dieser Quelle von \textsc{Bouvard} abgeleitete Resultat ist $\frac{1}{1070.5}$, und der wahrscheinliche Fehler dieser Bestimmung findet sich $= \pm \frac{1}{72\ObservedJupiterErrorMark}$ desselben, d. i. es ist gerade eben so wahrscheinlich, dass der wahre Werth zwischen $\frac{1}{1069}$ und $\frac{1}{1072}$ liegt, als dass er ausserhalb dieser Grenzen fällt. \textsc{Laplace} hat vorgezogen, die ungemein grosse Unwahrscheinlichkeit von etwas weitern Grenzen anzugeben: man kann 1 gegen fast eine Million wetten, dass der Fehler nicht mehr als $\frac{1}{100}$ des Ganzen beträgt. Hier ereignet sich nun der höchst merkwürdige Umstand, dass die von \textsc{Gauss} aus seiner Theorie der Pallasstörungen durch Jupiter abgeleitete Masse des Jupiter doch beträchtlich von jener Bestimmung verschieden ist, und weit ausserhalb jener Grenzen fällt, und was die Hauptsache ist, dass die Wahrscheinlichkeitstheorie auf diese Bestimmung angewandt, ihr sehr nahe eine eben so kleine Ungewissheit beilegt; über letzteres darf man sich nicht wundern, da bei den Bewegungen der Pallas die noch nicht so grosse Anzahl der Beobachtungen durch den weit stärkern Einfluss des Jupiter schon jetzt ersetzt wird. Wie soll man nun diesen Widerspruch ausgleichen, und welches ist die wahre Masse des Jupiter? Die fortgesetzten Beobachtungen der Pallas und die Untersuchung der Störungen der Juno werden in Zukunft hierüber weitere Auskunft geben. Erlaubt aber scheint ein Zweifel, ob die \textsc{Laplace}'sche Methode die Störungen zu berechnen, auf den Saturn angewandt, ganz so genaue Resultate gebe, als zu einer so delicaten Untersuchung erfordert werden. --- Ueber \textsc{Nonius} Formeln für die Dämmerung von \textsc{Delambre}. Das Verdienst von \textsc{Nonius} um die Aufgaben, die sich auf die Dauer der Dämmerung beziehen, war bisher nicht genug gewürdigt. Hr. \textsc{Delambre} zeigt in dieser sehr ausführlichen Abhandlung, dass in der That \textsc{Nonius} alle dahin gehörigen Aufgaben sehr gut und zum Theil besser als seine Nachfolger aufgelöst hat. Wir haben diese Abhandlung als eine Probe eines grossen Werks anzusehen, welches Hr. \textsc{Delambre} unter dem Titel einer Geschichte der Astronomie herausgeben will, und worin alle bekannten Bücher, wenigstens diejenigen, welche irgend etwas nützliches oder merkwürdiges enthalten, excerpirt, commentirt, zuweilen neu abgedruckt, aber immer wenigstens in neuer Einkleidung geliefert werden sollen. --- Der letzte Aufsatz dieses Jahrganges enthält die Elemente und Störungsgleichungen, nach welchen \textsc{Bouvard}'s neue Saturnstafeln berechnet sind, und die Vergleichung derselben mit allen seit 1747 beobachteten Oppositionen und Quadraturen dieses Planeten, zusammen 130. Wir hätten sehr gewünscht, dass alle diese den Tafeln zum Grunde liegenden beobachteten Oerter mit abgedruckt wären; so gut auch die Beobachtungen durch die Tafeln dargestellt werden, so wird man doch über kurz oder lang neue Verbesserungen zu suchen haben, wobei man auch auf alle jene Oppositionen wieder zurückkommen muss.
